% !TeX root = ../../Book3.tex
% 纲要标题：### 5.6*　限制根与对称空间分类
% 纲要位置：几何学计划纲要.md，第 2309 行。

% 纲要内容开始
% 
% > 写作范围与外部输入：规划任务：在第5.4节的对称分解与对偶模型之后建立限制根语言，完整计算一个经典矩阵商的秩、限制根和重数。给出不可约单连通紧型与非紧型模型的分类表及准确适用条件；一般 Cartan 分类作为外部输入，不以第一卷第19章的有限根算例冒充完整分类证明。所用根系与李群覆盖结果按实际版本就近恢复。
% 
% * 限制根、根空间与重数
% * 特殊线性群商的完整计算
% * 不可约经典模型与低秩同构
% * 例外模型及分类结论的外部输入
% * 局部等距、单连通覆盖与整体商
% * 分类表与紧非紧对偶的核验
% 
% ---
% 
% 纲要内容结束
% 正文状态：未撰写；本文件目前是框架。
% 规划任务：在第5.4节的对称分解与对偶模型之后建立限制根语言，完整计算一个经典矩阵商的秩、限制根和重数。给出不可约单连通紧型与非紧型模型的分类表及准确适用条件；一般 Cartan 分类作为外部输入，不以第一卷第19章的有限根算例冒充完整分类证明。所用根系与李群覆盖结果按实际版本就近恢复。

\OptionalSection{限制根与对称空间分类}{g3:sec:0506}

待撰写

\subsection{限制根、根空间与重数}
\label{g3:subsec:0506:01}
待撰写

\subsection{特殊线性群商的完整计算}
\label{g3:subsec:0506:02}
待撰写

\subsection{不可约经典模型与低秩同构}
\label{g3:subsec:0506:03}
待撰写

\subsection{例外模型及分类结论的外部输入}
\label{g3:subsec:0506:04}
待撰写

\subsection{局部等距、单连通覆盖与整体商}
\label{g3:subsec:0506:05}
待撰写

\subsection{分类表与紧非紧对偶的核验}
\label{g3:subsec:0506:06}
待撰写
